% === IX. CONCLUSION ==========================================================
\section{Conclusion and Future Work}
\label{sec:conclusion}

This paper studied the identifiability of a learning-based on-track Pacejka
identification pipeline~\cite{ontrack2025} at full vehicle scale, using a CARLA
simulation architecture that exposes the plant's own per-wheel forces as ground
truth. Transplanted unchanged from a 1:10 platform, the pipeline fails in a
specific, diagnosable and correctable way. The failure is a loss of excitation
created by the pipeline's own virtual-data step, which no change of solver,
bounds or tuning removes, and whose reachable friction
$\mu_{\mathrm{reach}}=v^{2}\delta_{\mathrm{sweep}}/(gL)$ depends on the rollout
configuration and not on the tire. A rollout speed appropriate at 1:10 caps
$\mu_{\mathrm{reach}}$ at \num{0.43} at full scale, below the plant's peak, so
the peak factor is not identifiable from the synthetic data and the bounded
optimiser resolves the ambiguity on a box edge.

Two further properties of the loop follow from the same analysis. The fit's
residual depends on $B$, $C$ and $D$ only through their product, because
stage~3 reconstructs its forces from the rollout's own quasi-steady yaw balance
and so observes nothing but the lateral acceleration the previous coefficient
set produced; the peak factor is therefore not recoverable from the rollout at
any excitation and must be measured on the recorded trace and held. And because
each fit is handed the trajectory its own previous answer generated, the loop
is a positive feedback path rather than a contraction, so stage~4 is
reformulated as a relaxed fixed-point iteration.

The correction comes with an extrapolation bound on the steering sweep, a
frozen regularisation target, identification against the achieved rather than
the commanded road-wheel angle, measured mass and geometry, a friction
warm start identified from the curvature of the axle force curve and pinned for
the loop, a relaxation $\beta=\num{0.20}$ on the coefficient update, and a
dynamics residual normalised to the units of the data term it is added to.
Together these produce identifications with no railed coefficient and the
correct understeer sign. Run against the inherited configuration over three
timed laps on a \num{0.70} surface, deliberately not the \num{1.05} the
configured prior was measured on, so that returning the prior scores badly,
the corrected pipeline recovers axle cornering stiffness to
\SI{3.2}{\percent} and \SI{1.4}{\percent} and the peak factor to
\SI{11}{\percent}, against an inherited configuration that rails $D$ on its
\num{2.0} box constraint in every cycle and overstates stiffness by
\SIrange{118}{132}{\percent}. Axle force RMSE is \num{12.1}$\times$ and
\num{10.5}$\times$ lower and MPC lateral tracking error \SI{32}{\percent}
lower, from \SI{0.224}{\meter} to \SI{0.151}{\meter} RMS, at unchanged solver
cost. Those numbers are an identification and not a preserved prior, with two
qualifications. The peak factor is identified by the brush estimator on the
recorded trace, not by the Pacejka fit, which Section~\ref{subsec:bcd} shows
cannot identify it; and one surface is not adaptivity, which needs a friction
\emph{change} the runs do not contain. That
the peak factor is recoverable from the rollout \emph{when the excitation
reaches it} is shown on a known tire in Section~\ref{subsec:sweep}, not on
these laps. Each configuration is a single
run, so the closed-loop margins are observations rather than estimates with a
confidence interval. We further argue that
per-coefficient box constraints are not a sufficient admissibility criterion
for handing a model to a model-based controller, and place gates on derived
physical quantities at that interface instead. The earlier reference
trajectory, which demanded \SI{1.22}{\g} from a vehicle sustaining
\SI{1.0}{\g}, was regenerated because that gate refused it rather than
loosened for it.

\subsection{Future work}

\paragraph{Immediate}
\begin{enumerate}
\item Run the ablations of Section~\ref{subsec:ablations} in the loop. The
relaxation sweep, the $D$-pinning comparison and the S4D and objective
ablations are offline replays of the identification stage on recorded buffers,
which is enough to attribute the result but not to confirm that the ordering
survives closed-loop coupling. Each is a scenario-file row.
\item Sweep the friction, rather than merely mismatching it. These runs put the
prior \SI{44}{\percent} away from a \emph{static} surface, which shows the
pipeline moves on evidence but not that it tracks. The step and decay schedules
already written in the scenario file are the intended vehicle, and are the
strongest available argument for identifying on track at all.
\item Excite the peak. The lap reaches
$P_{99}(|\alpha|)=\SI{2.7}{\degree}$ and \SI{5.8}{\degree} at most, against a
plant whose curve peaks at \SI{8.0}{\degree} on this surface, and an axle
utilisation of \num{0.71}. That is enough for the friction gate to release,
and it is also where the brush estimator's \SI{11}{\percent} bias comes from. A
trajectory or an injected manoeuvre that traverses the peak would remove it.
\item Run the noise sweep of~\cite{ontrack2025} on this platform using the
bridge's odometry noise model, with a nonlinear least-squares identification of
the same data as the comparison arm.
\item Cross-check the friction warm start against the simulator's per-wheel
forces across a pace sweep, including the $\sigma_{\mu}/\mu$ gate boundary and
the IMU-versus-kinematics cross-check that guards the mass and inertia
assumptions. It is exercised at one excitation level on the CARLA runs, where
it returns $\mu$ \SI{11}{\percent} high, and swept only on a synthetic
brush-tire plant where $\mu$ is known exactly (Table~\ref{tab:mu}). Since the
published peak factor now \emph{is} that estimate, its bias is the pipeline's
bias.
\item Report lap time on the simulator's own clock. Lap timing currently runs
on the ROS wall clock against a faster-than-real-time simulator, which makes
lap time unusable as a comparison metric (Section~\ref{subsec:resultsclosedloop}).
\item Measure $I_{z}$ rather than assuming it
(Section~\ref{subsec:assumptions}).
\end{enumerate}

\paragraph{Beyond}
\begin{enumerate}
\item \textbf{Excitation-aware data collection.} Equation~\eqref{eq:mureach}
inverts into a requirement on the manoeuvre. A controller that deliberately
inserts bounded slalom or ramp-steer excitation when the identifier reports
insufficient slip-angle coverage would move the pipeline from passive on-track
identification toward active, safety-bounded on-track experiment design. That
is optimal input design~\cite{goodwin1977design} under a safety constraint,
and is where the literature for this step already exists.
\item \textbf{Uncertainty-aware output.} The gates of
Section~\ref{subsec:gates} are binary. Reporting a posterior over the
coefficients, through the Bayesian variational path already implemented
alongside the deterministic solvers, would let the controller weight the model
by its confidence rather than accept or reject it.
\item \textbf{Hardware replication.} The decisive test of every claim in
Section~\ref{sec:disc} is a physical Formula Student vehicle. The
identification code is unchanged between simulation and hardware. What changes
is the availability of ground truth, and that is what makes the simulation
study a necessary precursor rather than a substitute.
\end{enumerate}
